% Einheit 3 von 4 – Antiproportionale Zuordnungen (bank/zuordnungen/e3.jsonl, zusammenbau v0.1)
%% TODO Zweigzeile Teil 1: was hier gelernt wird (Satz in Schülersprache aus dem Einheitstitel) – Einheit 3
\einheitenkopf[e3]{Einheit 3 von 4 · Antiproportionale Zuordnungen}
\zweigzeile{ab Kl. 7 · P10}
\setcounter{aufgabe}{19}

%% TODO Titel: Ich-kann-Satz fehlt in der Bank (Platzhalter: Kettenname) – Nr. 20
%% TODO Anweisung: ein Satz über der ersten Teilaufgabe fehlt in der Bank – Nr. 20
\begin{aufgabe}{Antiproportional}
%% TODO \rechenplatz{n}: Schrittzahl des Grundfalls fehlt in der Bank (nur bei fester Schreibform, 2.3 b)
\begin{teile}
\teil Je mehr Maler an einer Wand arbeiten, desto … Zeit brauchen sie. \\ \kreuz{mehr} \\ \kreuz{weniger}
\teil 2 Bagger brauchen für eine Baugrube 10 Tage. Wie lange brauchen 4 Bagger? \leerfeld[Tage]
\teil 2 Traktoren pflügen ein Feld in 9 Stunden. Wie lange brauchen 6 Traktoren? \leerfeld[h]
\teil 3 Maschinen brauchen für einen Auftrag 8 Stunden. Wie lange brauchen 4 Maschinen? \leerfeld[h]
\teil Jemand behauptet: 6 Arbeiter brauchen 10 Tage, 4 Arbeiter brauchen 15 Tage. Prüfe mit dem Produkt, ob das zusammenpasst.
\teil Für einen Umzug haben 3 Freunde 4 Stunden eingeplant. Kurzfristig helfen 3 weitere Freunde mit. Alle arbeiten gleich schnell. Wie lange dauert der Umzug jetzt?
\teil Ein Futtervorrat reicht für 6 Ziegen 10 Tage. Wie viele Tage reicht er für 5 Ziegen? \hfill (P10 2014 GYM) \\ \leerfeld[Tage]
\end{teile}
\end{aufgabe}

%% TODO Titel: Ich-kann-Satz fehlt in der Bank (Platzhalter: Kettenname) – Nr. 21
%% TODO Anweisung: ein Satz über der ersten Teilaufgabe fehlt in der Bank – Nr. 21
\begin{aufgabe}{Graph als fallende Kurve, Werte ablesen}
\begin{teile}
\teil Lies ab: Wie lange brauchen 3 Helfer? Wie viele Helfer braucht man, damit es 2 Stunden dauert?

\begin{ksys}[xmin=0,xmax=7,ymin=0,ymax=13,ablesen,xlabel=Anzahl Helfer,ylabel=Zeit in h]\funktionab{12/\x}{}{1}{6}\end{ksys}
\leerfeld[h] ; \leerfeld Helfer
\end{teile}
\end{aufgabe}

%% TODO Titel: Ich-kann-Satz fehlt in der Bank (Platzhalter: Kettenname) – Nr. 22
%% TODO Anweisung: ein Satz über der ersten Teilaufgabe fehlt in der Bank – Nr. 22
\begin{aufgabe}{Antiproportional – Fehler finden, Begründen, Darstellungswechsel, Anwendung}
\begin{teile}
\teil 4 Maler brauchen 5 Tage. Kim rechnet: „8 Maler brauchen 10 Tage, denn es sind doppelt so viele.“ Finde den Fehler und rechne richtig.
\teil Warum bleibt bei der Zuordnung Anzahl Arbeiter → Anzahl Tage das Produkt gleich? Begründe.
\teil Trage die Punkte der Tabelle ein und verbinde sie zu einer Kurve.

\wertetabelle[6,3,2,1]{Anzahl Kinder}{Stück je Kind}{1,2,3,6} \begin{ksys}[xmin=0,xmax=7,ymin=0,ymax=7,xlabel=Anzahl Kinder,ylabel=Stück je Kind]\end{ksys}
\teil Für das Schmücken der Turnhalle plant die Klasse mit 5 Helfern 3 Stunden. Die Halle ist aber nur 2 Stunden frei. Wie viele Helfer braucht man mindestens?
\end{teile}
\end{aufgabe}
